\documentclass[10pt,a4paper]{amsart}
\usepackage[margin=19mm]{geometry}
\usepackage{amsmath,amssymb,mathtools}
\usepackage[scr=rsfs,cal=euler]{mathalfa}
\usepackage{xfrac}
\usepackage[unicode,hidelinks,bookmarks=false]{hyperref}
\newcommand{\ie}{i.e.\ }
\newcommand{\cf}{cf.\ }
\newcommand{\resp}{respectively\ }
\newcommand{\Subsection}{\subsection}
\providecommand{\nobreakdash}{\nobreak\hbox{-}}
\setlength{\emergencystretch}{2em}
\multlinegap0pt
\usepackage{mathtools}
\usepackage{amssymb}
\newtheorem{proposition}{Proposition}
\title{Heat semigroups on quantum automorphism groups of finite dimensional C*-algebras}
\author{Futaba Sato}
\date{Source: arXiv:2503.03448v4 (CC BY 4.0)}
\begin{document}
\maketitle

\noindent\emph{Source and licence.} Heat semigroups on quantum automorphism groups of finite dimensional C*-algebras. Version arXiv:2503.03448v4 (CC BY 4.0), distributed by arXiv under the Creative Commons Attribution 4.0 International licence, http://creativecommons.org/licenses/by/4.0/, as recorded in the arXiv licence metadata for this exact version and shown on the abstract page https://arxiv.org/abs/2503.03448v4. Full licence notice: \texttt{licenses/arxiv-2503.03448-CC-BY-4.0.txt}.\par\smallskip

\noindent\textit{Editorial scope.} Counted unit: the article's proposition proposition (source0.tex lines 385--389). The article's complete original proof of it is delivered with it (source0.tex lines 391--419, one \texttt{proof} environment of 29 lines). No mathematics is written, repaired or re-derived: the statement and the proof are the article's own lines, in the article's own notation, and no retained line is substituted or re-flowed.  No further source-local context is needed: every cross-reference of every retained line resolves inside the two delivered spans.  Nothing was omitted from any retained span, so the package carries no omission record.  No retained line contains a citation, so the wrapper carries no bibliography and no bibliography entry.  No retained line contains a figure environment, an image-inclusion command or drawing code, so no image is delivered and the package declares no asset dependency.  Editorial formatting only, in addition: an AMS \texttt{amsart} preamble that restates the article's own declarations and macros used by the retained lines, namely the environments \texttt{proposition} on separate counters, together with the standard packages those lines need.  The retained environments print in this document as Proposition 1 in order of appearance, whereas the article prints the same items with the numbering of its own full document.  The complete native source of the article as distributed in the arXiv e-print for this exact version is bundled unchanged as \texttt{sources/174256/source0.tex}.
 \begin{proposition}
 For a heat semigroup $\{T_t\}$ of ${\rm Aut}^+(B)$ with $\dim B\geq5$, there exists $\varepsilon_0\geq0$ such that for any $p\geq4-\varepsilon_0$,
 \[||T_t||_{2\to p}\leq1, \text{for\;all}\;t\geq\frac{dn}{2}\log{(p-1)}+(1-\frac{2}{p}n\log D)\]
 with $d=\frac{\log(11+\sqrt{105})-\log2}{\log3}.$
 \end{proposition}

 \begin{proof}
 By the H\"older inequality, for $p\geq1$, we have
 \[||x_k||_p\leq(D(k+1))^{1-\frac{2}{p}}||x_k||_2, \;x_k\in V_k.\]
 Therefore
 \[||T_t(x)||_p^2\leq |h(x)|^2+(p-1)\left(\sum_{k\geq1} e^{\lambda_k t}||x_k||_p\right)^2\]
 \[\leq |h(x)|^2+(p-1)\left(\sum_{k\geq1} e^{\lambda_k t}(D(2k+1))^{1-\frac{2}{p}}||x_k||_2\right)^2\]
 \[\leq  |h(x)|^2+(p-1) \sum_{k\geq1} e^{2 \lambda_k t}(D(2k+1))^{2(1-\frac{2}{p})}||x_k||_2^2.\]
 When $t\geq \frac{dn}{2}\log{(p-1)}+(1-\frac{2}{p})n\log{D}$ and $s\geq1$, 
 \[2\lambda_kt\leq-dk\log{(p-1)}-2\left(1-\frac{2}{p}\right)k\log{D}\]
 \[\leq -dk\log{(p-1)}-2\left(1-\frac{2}{p}\right)\log{D}\]
 and $e^{2\lambda_kt}\leq (p-1)^{-dk}D^{-2(1-\frac{2}{p})}.$
 
Consider $\phi(p)\coloneqq (p-1)^{1-dk}(2k+1)^{2(1-\frac{2}{p})}$.
 Therefore it suffices to show that for some $\varepsilon_0$, for any $p\geq 4-\varepsilon_0$,
 \[R_p\coloneqq\sum_{k\geq1}\phi(p)=\sum_{k\geq1}(p-1)^{1-dk}(2k+1)^{2(1-\frac{2}{p})}\leq 1.\]
Note that 
 \[R_4=\sum_{k\geq1}\frac{2k+1}{3^{dk-1}}=\frac{3(3\cdot3^d-1)}{(3^d-1)^2}\]
 and $d$ needs to satisfy $d\geq \frac{\log(11+\sqrt{105})-\log2}{\log3}$.
 
We have that $\phi'(p)\leq0$ if and only if
\[\frac{4(p-1)}{p^2}\leq\frac{dk-1}{\log{(2k+1)}}.\]
As $f_1(p)=\frac{4(p-1)}{p^2}$ 
is decreasing for $p\geq2$ and $f_2(k)=\frac{dk-1}{\log{(2k+1)}}$
 is increasing for 
 $k\geq1$, $d=\frac{\log(11+\sqrt{105})-\log2}{\log3}>\log{3}+1$ satisfies the following:
\[f_1(p)\leq f_1(2)=1<\frac{d-1}{\log{3}}=f_2(1)\leq f_2(s)\]
for any $p\geq2, s\geq1$.
Now we have that $d =\frac{\log(11+\sqrt{105})-\log2}{\log3}$ satisfies the desired condition.
 \end{proof}

\end{document}
